\section{Feasibility Study}
\label{sec:visibility}
Exploiting DPSs' acoustic side-channel vulnerability to recover speech signals requires a thorough understanding of the sensor's behavior, sampling rate, and frequency response. 
Sensor datasheets often conceal these details, which presents challenges for attackers attempting to reconstruct audio from pressure readings. Our analysis aims to highlight these challenges and demonstrate the attack's feasibility.

\begin{figure}[t!]
\centering
\includegraphics[trim={1px 1px 1px 1px},clip,width=0.75\linewidth]{newestPDF/Expt_setup.pdf}
\vspace{-0.5em}
  \caption{Our experimental setup for a feasibility study.}
  \label{fig:experiment}
  \vspace{-0.8em}
\end{figure}


\subsection{Experimental setup}
\label{subsection:experimental_setup}
The experimental setup is depicted in Fig. \ref{fig:experiment}. To showcase our attack's feasibility, we use the SDP800 DPS~\cite{datasheetsdp} (see~\ref{subsection:physicsofDPS}), securely mounted to avoid movement-related noise. 
A Sony XB13 \cite{xb13} speaker is placed 5 cm directly underneath one of the pressure ports to play the audio clips. 
To prevent sound signals from influencing the pressure in the other port, we connect a sampling tube to it, isolated by positioning its opening a meter away, shielded by a hard surface. 
Signals from the sensor are read using a Raspberry Pi 3 Model B~\cite{rs-online}.
We employ a custom Python script to convert the pressure readings into a format an attacker can process. 

\begin{figure}[b]
% \vspace{-00.120em}
\centering
\includegraphics[trim={10px 22px 10px 18px},clip,width=0.49\textwidth]{newPDF/one_STFT.pdf}
  \vspace{-01.60em}
  \caption{STFT plot of a spoken word "one".}
  \label{fig:one_stft}
\vspace{-0.800em}
\end{figure}


\subsection{Primary observation}
Fig.~\ref{fig:one_stft} compares the Short-Time Fourier Transform (STFT) spectrum of the word "one" captured by an iPhone 13 microphone (right) and the SDP800 DPS (left). 
The spectrum demonstrates the variations in power across different audio frequency components over time.
Notably, the DPS's response is discernibly feebler and less extensive than the microphone’s, especially in frequencies surpassing $0.4$ kHz. 
Distortions are observable even in zones of ostensibly strong signals ($0.4$ - $0.85$ kHz). 
Furthermore, the DPS fails to register signals exceeding $0.9$ kHz. 
A comparative analysis of the Fast Fourier Transform (FFT) plots in Fig.~\ref{fig:one_FFT} echoes these observations, with the DPS showing subpar responsiveness to signals above $0.4$ kHz.
{This indicates a notable loss in acoustic fidelity when reconstructing speech from the DPS data (also see Fig.~\ref{fig:nonLinearity}).} 

\begin{figure}[h]
% \vspace{-01.120em}
\centering
\includegraphics[trim={5px 20px 5px 20px},clip,width=0.41\textwidth]{newPDF/one_FFT_2.pdf}
\vspace{-0.5em}
  \caption{FFT plot of a spoken word "one".}
  \label{fig:one_FFT}
\vspace{-01.0em}
\end{figure}

\subsection{Challenges}
\label{subsection:challenges}
A successful attack requires recovery of information lost from pressure readings, at least partially.
Here, we examine the challenges attackers may face during this process.

\subsubsection{\textbf{Low sampling rate}} 
The sensor's low sampling rate is a significant challenge in executing the BaroVox attack.
For example, the SDP800 sensor used in our experiment has a theoretical maximum sampling rate of 2.2 kHz~\cite{datasheetsdp}.
However, as demonstrated in Fig.~\ref{fig:one_FFT}, it only captures frequencies up to approximately 0.9 kHz, with notably weak signals within the 0.6 - 0.9 kHz range. 
Consequently, given Nyquist’s theorem~\cite{vaidyanathan2001generalizations}, the sensor's effective sampling rate is approximated at 1.8 kHz, a shortfall of 0.4 kHz from its potential. 
This limitation restricts speech recovery within the sub-0.9 kHz range, significantly below the 4 kHz threshold necessary for intelligible speech~\cite{heide1998speech, villchur1973signal}, causing aliasing and complicating the recovery stage.   






\subsubsection{\textbf{Non-linear frequency response}}
\label{subsection:non-linearity}

DPS can be depicted as a second-order dynamic system since it employs an elastic diaphragm for pressure force collection~\cite{whitmore2009improved}.
This leads to a non-linear frequency response to sound waves of varying frequencies.
To investigate this, we construct an audio file with a sine sweep wave ranging from 1 Hz to 2 kHz. We play the audio file through a speaker using the setup specified in Sec. \ref{subsection:experimental_setup}. We then reconstruct audio from pressure reading and analyze the data using FFT plots.
FFT plots of the reconstructed audio (Fig. ~\ref{fig:nonLinearity}) show reduced signal strength across frequencies compared to the original. 
{The original audio has consistent power up to 2 kHz, whereas the reconstructed signal loses power as frequency increases, dropping sharply after 0.4 kHz. }
At frequencies above 0.65 kHz, its power nears mere noise levels.
These insights can be leveraged to equalize the signal by boosting the strength of the affected frequencies (see Sec.\ref{subsubsection:equalization}).

\subsubsection{\textbf{Low Signal-to-Noise Ratio (SNR)}}
\label{subsection:lowSNR} 
DPSs yield non-zero readings even in the absence of sound due to the presence of ambient noise and pressure fluctuations.
Acoustic disturbances like machinery operation, air conditioning noise, and door activities further affect the readings.
These ambient noises introduce unwanted interference to the sensor's data, lowering the SNR and degrading the quality of reconstructed speech.
To mitigate this issue, we applied filters and spectral subtraction techniques to improve the audio recovery quality (see Sec.~\ref{subsection:DS-1}).
